\chapter{Динамическое программирование}
\label{ch:dp}

\begin{quote}
\itshape
Дайте мне точку опоры, и я переверну Землю.
\upshape --- Архимед
\end{quote}

\bigskip

Динамическое программирование\index{dynamic programming} (DP) занимает уникальное место в обучении с подкреплением: оно решает центральную задачу последовательного принятия решений \emph{точно}, но лишь тогда, когда агент располагает идеальной математической моделью среды.  В этой главе мы рассматриваем DP как собрание практических алгоритмов и как теоретическую основу, на которой зиждется вся остальная книга.  Каждый алгоритм, с которым мы встретимся в последующих главах\,---\,методы Монте-Карло, временное разностное обучение, градиент политики\,---\,можно понять как приближение к DP, заменяющее точные модельные вычисления выборкой или аппроксимацией функций.

Формально мы предполагаем на протяжении всей этой главы, что среда является конечным\index{finite MDP} марковским процессом принятия решений $({\cS}, {\cA}, p, \gamma)$, как определено в главе~\ref{ch:mdp}.  Ядро переходов $p(s', r \mid s, a)$ известно полностью, пространства состояний и действий конечны, а коэффициент дисконтирования удовлетворяет $\gamma \in [0,1)$.  В этих условиях DP находит оптимальную функцию ценности и оптимальную политику за время, полиномиальное относительно $|\cS|$ и $|\cA|$.

Фундаментальное понимание, лежащее в основе всех алгоритмов DP, — \emph{уравнение Беллмана}\index{Bellman equation}: оптимальные функции ценности удовлетворяют рекуррентному условию самосогласованности, и DP эксплуатирует эту структуру, итерируя тщательно определённый оператор до сходимости.  Мы представляем четыре основных алгоритма\,---\,итеративную оценку политики, улучшение политики, итерацию политики и итерацию ценности\,---\,и объединяем их под знаменем \emph{обобщённой итерации политики} (GPI)\index{Generalized Policy Iteration}.

% ===========================================================
\section{Оценка политики (предсказание)}
\label{sec:policy-evaluation}
\index{policy evaluation}

Первая и наиболее базовая задача — \emph{оценка политики}\index{policy evaluation}: дана фиксированная политика $\pi$, вычислить её функцию ценности состояния $v^\pi$.  Это задача предсказания\,---\,мы ещё не пытаемся найти \emph{хорошую} политику, лишь понять, насколько хорошо данная.

\subsection*{Оператор Беллмана для фиксированной политики}

Из главы~\ref{ch:mdp} напомним, что $v^\pi$ является единственным решением \emph{уравнений ожиданий Беллмана}\index{Bellman expectation equation}
\begin{equation}
\label{eq:bellman-expectation}
v^\pi(s) = \sum_{a \in \cA} \pi(a \mid s) \sum_{s' \in \cS} \sum_{r} p(s', r \mid s, a)\bigl[r + \gamma\, v^\pi(s')\bigr], \qquad s \in \cS.
\end{equation}
Это линейная система из $|\cS|$ уравнений с $|\cS|$ неизвестными.  Теоретически её можно было бы решить напрямую методом Гаусса, но для больших пространств состояний, типичных в RL, это вычислительно нецелесообразно.  Вместо этого DP решает~\eqref{eq:bellman-expectation} \emph{методом последовательных приближений}: определим \textbf{оператор Беллмана для политики}\index{Bellman operator}
\begin{equation}
\label{eq:bellman-operator}
(T^\pi v)(s) = \sum_{a \in \cA} \pi(a \mid s) \sum_{s' \in \cS} \sum_{r} p(s', r \mid s, a)\bigl[r + \gamma\, v(s')\bigr],
\end{equation}
и итерируем $v_{k+1} = T^\pi v_k$ начиная с произвольной ограниченной инициализации $v_0$.

\begin{theorem}[Сжатие $T^\pi$]
\label{thm:contraction}
\index{contraction mapping}
Оператор $T^\pi$ является сжатием\index{Banach fixed-point theorem} на пространстве ограниченных функций $\cS \to \R$, оснащённом супремум-нормой $\|v\|_\infty = \max_{s \in \cS}|v(s)|$.  Точнее, для любых двух ограниченных функций $u, v \colon \cS \to \R$,
\[
\|T^\pi u - T^\pi v\|_\infty \le \gamma \|u - v\|_\infty.
\]
Следовательно, $T^\pi$ имеет единственную неподвижную точку $v^\pi$, и итераты $v_{k+1} = T^\pi v_k$ сходятся к $v^\pi$ из любой начальной точки $v_0$.
\end{theorem}

\begin{proof}
Зафиксируем произвольное состояние $s \in \cS$.  Тогда
\begin{align*}
\bigl|(T^\pi u)(s) - (T^\pi v)(s)\bigr|
&= \Bigl|\sum_a \pi(a \mid s) \sum_{s',r} p(s',r \mid s,a)\,\gamma\bigl[u(s') - v(s')\bigr]\Bigr| \\
&\le \sum_a \pi(a \mid s) \sum_{s',r} p(s',r \mid s,a)\,\gamma\,\|u - v\|_\infty \\
&= \gamma \|u - v\|_\infty,
\end{align*}
где неравенство использует неравенство треугольника и тот факт, что веса $\pi(a \mid s)\,p(s',r \mid s,a)$ неотрицательны и суммируются к единице.  Взяв супремум по $s$, получаем $\|T^\pi u - T^\pi v\|_\infty \le \gamma\|u-v\|_\infty$.  Поскольку $\gamma < 1$, оператор является строгим сжатием.  По теореме Банаха о неподвижной точке он имеет единственную неподвижную точку, а итераты сходятся с геометрической скоростью с коэффициентом $\gamma$.
\end{proof}

\begin{remark}
После $k$ итераций ошибка удовлетворяет $\|v_k - v^\pi\|_\infty \le \gamma^k \|v_0 - v^\pi\|_\infty$.  Для $\gamma = 0.9$ и желаемой точности $\varepsilon = 10^{-3}$ требуется примерно $k \ge \log(10^{-3}/\|v_0 - v^\pi\|_\infty)/\log(0.9)$ итераций\,---\,независимо от $|\cS|$.
\end{remark}

\subsection*{Алгоритм и критерий остановки}

На практике обходят все состояния и применяют обновление одновременно (синхронное DP\index{synchronous DP}) или последовательно (см. раздел~\ref{sec:async-dp}).  Алгоритм останавливается, когда максимальное изменение по всем состояниям падает ниже порога $\theta$.

\begin{algorithm}[ht]
\caption{Итеративная оценка политики}
\label{alg:policy-eval}
\KwIn{Политика $\pi$, динамика MDP $p$, дисконтирование $\gamma$, порог $\theta > 0$}
Инициализировать $V(s) \leftarrow 0$ для всех $s \in \cS$; установить $V(\text{терминал}) \leftarrow 0$\;
\Repeat{$\Delta < \theta$}{
  $\Delta \leftarrow 0$\;
  \For{каждого $s \in \cS$}{
    $v \leftarrow V(s)$\;
    $V(s) \leftarrow \displaystyle\sum_a \pi(a \mid s) \sum_{s',r} p(s',r \mid s,a)\bigl[r + \gamma\, V(s')\bigr]$\;
    $\Delta \leftarrow \max(\Delta,\; |v - V(s)|)$\;
  }
}
\KwOut{$V \approx v^\pi$}
\end{algorithm}

\paragraph{Критерий остановки.}\index{stopping criterion}
Порог $\theta$ управляет компромиссом между точностью и вычислительными затратами.  Распространённый практический выбор — $\theta = 10^{-4}$ или $\theta = 10^{-6}$.  \emph{А-posteriori} оценка ошибки равна
\[
\|V - v^\pi\|_\infty \le \frac{\theta\,\gamma}{1-\gamma},
\]
что следует из аргумента геометрического ряда, применённого к сжатию.

\subsection*{Числовой пример: решётка}

\begin{example}[$4\times4$ решётка]
\label{ex:gridworld-eval}
\index{gridworld}
Рассмотрим стандартную $4\times4$ решётку из 16 состояний, расположенных в сетке.  Состояния $s_1$ (верхний левый) и $s_{16}$ (нижний правый) — терминальные.  Из каждого нетерминального состояния агент может двигаться \textsc{вверх}, \textsc{вниз}, \textsc{влево} или \textsc{вправо}; попытка выйти за пределы решётки оставляет состояние неизменным.  Каждый переход даёт вознаграждение $-1$ (агент хочет достичь терминального состояния как можно быстрее).  Оцениваем \emph{равновероятную случайную политику} $\pi(a \mid s) = 0.25$ для всех $a$ и $\gamma = 1$ (недисконтированная эпизодическая задача).

\medskip
\noindent
В следующей таблице показаны значения $V$ после нескольких проходов.

\begin{center}
\begin{tabular}{ccccc}
\toprule
Итерация & \multicolumn{4}{c}{Значения $V$ (строка за строкой, терминальные состояния показаны как $0$)} \\
\midrule
$k=0$ & $0$ & $0$ & $0$ & $0$ \\
      & $0$ & $0$ & $0$ & $0$ \\
      & $0$ & $0$ & $0$ & $0$ \\
      & $0$ & $0$ & $0$ & $0$ \\
\midrule
$k=1$ & $0$ & $-1$ & $-1$ & $-1$ \\
      & $-1$ & $-1$ & $-1$ & $-1$ \\
      & $-1$ & $-1$ & $-1$ & $-1$ \\
      & $-1$ & $-1$ & $-1$ & $0$ \\
\midrule
$k=2$ & $0$ & $-1.75$ & $-2$ & $-2$ \\
      & $-1.75$ & $-2$ & $-2$ & $-2$ \\
      & $-2$ & $-2$ & $-2$ & $-1.75$ \\
      & $-2$ & $-2$ & $-1.75$ & $0$ \\
\midrule
$k \to \infty$ & $0$ & $-14$ & $-20$ & $-22$ \\
               & $-14$ & $-18$ & $-20$ & $-20$ \\
               & $-20$ & $-20$ & $-18$ & $-14$ \\
               & $-22$ & $-20$ & $-14$ & $0$ \\
\bottomrule
\end{tabular}
\end{center}

\medskip
\noindent
Сошедшиеся значения отражают ожидаемое число шагов до достижения терминального состояния под случайной политикой.  Состояние, диагонально противоположное обоим терминалам ($-22$), труднее всего покинуть; состояния, соседние с терминалами, имеют $v^\pi \approx -14$.  Этот пример иллюстрирует, как оценка политики распространяет информацию \emph{обратно} от терминальных состояний через решётку.
\end{example}

\begin{figure}[ht]
\centering
\begin{tikzpicture}[>=Stealth, node distance=0.9cm,
  cell/.style={draw, minimum size=0.85cm, font=\scriptsize},
  tcell/.style={draw, minimum size=0.85cm, font=\scriptsize, fill=gray!30}]
% Рисуем сетку 4x4 со сошедшимися значениями
\foreach \row/\vals in {
  0/{0, -14, -20, -22},
  1/{-14, -18, -20, -20},
  2/{-20, -20, -18, -14},
  3/{-22, -20, -14, 0}}{
  \foreach \col/\val [count=\ci from 0] in \vals {
    \pgfmathsetmacro{\isterm}{(\row==0 && \ci==0) || (\row==3 && \ci==3) ? 1 : 0}
    \ifnum\ci=0
      \def\vval{\row}% trick
    \fi
  }
}
% Ручная сетка для ясности
\foreach \r in {0,1,2,3}{
  \foreach \c in {0,1,2,3}{
    \pgfmathtruncatemacro{\istop}{(\r==0 && \c==0) || (\r==3 && \c==3)}
    \ifnum\istop=1
      \node[tcell] at (\c*0.95, -\r*0.95) {$0$};
    \else
      % вычислить значение
      \pgfmathtruncatemacro{\idx}{\r*4+\c}
      \pgfmathsetmacro{\vval}{
        \idx==1 ? -14 : (\idx==2 ? -20 : (\idx==3 ? -22 :
        (\idx==4 ? -14 : (\idx==5 ? -18 : (\idx==6 ? -20 :
        (\idx==7 ? -20 : (\idx==8 ? -20 : (\idx==9 ? -20 :
        (\idx==10 ? -18 : (\idx==11 ? -14 : (\idx==12 ? -22 :
        (\idx==13 ? -20 : (\idx==14 ? -14 : 0)))))))))))))))}
      \node[cell] at (\c*0.95, -\r*0.95) {\pgfmathprintnumber[fixed,precision=0]{\vval}};
    \fi
  }
}
\node[above] at (1.425, 0.55) {\small Сошедшаяся $v^\pi$ (случайная политика, $\gamma=1$)};
\end{tikzpicture}
\caption{Функция ценности состояния $v^\pi$ для равновероятной случайной политики на $4\times 4$ решётке после сходимости итеративной оценки политики.  Закрашенные клетки — терминальные состояния со значением $0$.  Значения представляют ожидаемое суммарное (недисконтированное) вознаграждение, то есть отрицательное ожидаемое число шагов до достижения терминала.}
\label{fig:gridworld-values}
\end{figure}

% ===========================================================
\section{Улучшение политики}
\label{sec:policy-improvement}
\index{policy improvement}

Вычислив $v^\pi$ для некоторой политики $\pi$, естественно спросить: можно ли построить строго лучшую политику?  \emph{Теорема об улучшении политики}\index{policy improvement theorem} отвечает на это утвердительно.

\subsection*{Функция ценности действия}

Дано $v^\pi$, \emph{функция ценности действия}\index{action-value function} $q^\pi(s,a)$ фиксирует ценность выполнения действия $a$ в состоянии $s$ с последующим следованием $\pi$:
\begin{equation}
\label{eq:q-from-v}
q^\pi(s,a) = \sum_{s' \in \cS}\sum_r p(s',r \mid s,a)\bigl[r + \gamma\,v^\pi(s')\bigr].
\end{equation}
В модельной постановке $q^\pi$ можно вычислить непосредственно из $v^\pi$ с помощью одношагового просмотра вперёд в~\eqref{eq:q-from-v}.  Ключевое наблюдение: $q^\pi(s, a) \ge v^\pi(s)$ означает, что выполнение действия $a$ \emph{один раз} в состоянии $s$ с последующим возвратом к $\pi$ не хуже, чем следование $\pi$ с самого начала.

\begin{theorem}[Теорема об улучшении политики]
\label{thm:policy-improvement}
\index{policy improvement theorem}
Пусть $\pi$ и $\pi'$ — две детерминированные политики такие, что для всех $s \in \cS$,
\begin{equation}
\label{eq:improvement-condition}
q^\pi\bigl(s, \pi'(s)\bigr) \ge v^\pi(s).
\end{equation}
Тогда $\pi'$ не хуже $\pi$:
\[
v^{\pi'}(s) \ge v^\pi(s), \qquad \forall s \in \cS.
\]
Более того, если неравенство в~\eqref{eq:improvement-condition} строгое хотя бы для одного состояния, то $v^{\pi'}(s) > v^\pi(s)$ хотя бы для одного состояния.
\end{theorem}

\begin{proof}
Из~\eqref{eq:improvement-condition} и определения $q^\pi$:
\[
v^\pi(s) \le q^\pi\bigl(s, \pi'(s)\bigr)
= \sum_{s',r} p(s',r \mid s, \pi'(s))\bigl[r + \gamma\,v^\pi(s')\bigr]
= (T^{\pi'} v^\pi)(s).
\]
Следовательно, $v^\pi \le T^{\pi'} v^\pi$ поточечно.  Оператор $T^{\pi'}$ \emph{монотонен}: если $u \le v$ поточечно, то $(T^{\pi'} u)(s) \le (T^{\pi'} v)(s)$ для всех $s$, поскольку веса $p(s',r \mid s, \pi'(s))$ неотрицательны.  Применяя $T^{\pi'}$ многократно и используя монотонность,
\[
v^\pi \le T^{\pi'} v^\pi \le (T^{\pi'})^2 v^\pi \le \cdots \le (T^{\pi'})^n v^\pi \qquad \forall n \ge 1.
\]
По теореме~\ref{thm:contraction} (с $\pi'$ вместо $\pi$), $(T^{\pi'})^n v^\pi \to v^{\pi'}$ при $n \to \infty$.  Поскольку неравенство сохраняется при каждом конечном $n$, оно сохраняется в пределе:
\[
v^\pi(s) \le v^{\pi'}(s), \qquad \forall s \in \cS. \qedhere
\]
\end{proof}

\subsection*{Построение жадной политики}

\emph{Жадная политика}\index{greedy policy} $\pi'$ относительно $v^\pi$ удовлетворяет условию улучшения~\eqref{eq:improvement-condition} максимальным образом:
\begin{equation}
\label{eq:greedy-policy}
\pi'(s) \in \argmax_{a \in \cA} q^\pi(s,a)
= \argmax_{a \in \cA} \sum_{s',r} p(s',r \mid s,a)\bigl[r + \gamma\,v^\pi(s')\bigr].
\end{equation}
По построению, $q^\pi(s, \pi'(s)) = \max_a q^\pi(s,a) \ge v^\pi(s)$ для всех $s$, поэтому теорема применима.

\paragraph{Когда улучшение прекращается.}\index{optimality condition}
Предположим, жадное улучшение даёт $\pi' = \pi$\,---\,то есть текущая политика уже жадна относительно своей собственной функции ценности.  Тогда $v^\pi(s) = \max_a q^\pi(s,a)$ для всех $s$, что в точности является \emph{уравнением оптимальности Беллмана}\index{Bellman optimality equation}:
\[
v^\pi(s) = \max_{a \in \cA} \sum_{s',r} p(s',r \mid s,a)\bigl[r + \gamma\, v^\pi(s')\bigr] = v^*(s).
\]
Следовательно, $\pi$ является оптимальной политикой.  Это наблюдение движет алгоритмом итерации политики.

% ===========================================================
\section{Итерация политики}
\label{sec:policy-iteration}
\index{policy iteration}

Итерация политики чередует два шага, разработанных выше: оцениваем текущую политику до сходимости, затем жадно улучшаем её.

\begin{algorithm}[ht]
\caption{Итерация политики}
\label{alg:policy-iteration}
\KwIn{Динамика MDP $p$, дисконтирование $\gamma$, порог $\theta > 0$}
Произвольно инициализировать $V(s) \in \R$ и $\pi(s) \in \cA$ для всех $s \in \cS$; $V(\text{терминал}) \leftarrow 0$\;
\Repeat{$\textsc{политика-стабильна} = \mathtt{true}$}{
  \tcp{Оценка политики}
  \Repeat{$\Delta < \theta$}{
    $\Delta \leftarrow 0$\;
    \For{каждого $s \in \cS$}{
      $v \leftarrow V(s)$\;
      $V(s) \leftarrow \displaystyle\sum_{s',r} p(s',r \mid s,\pi(s))\bigl[r + \gamma\,V(s')\bigr]$\;
      $\Delta \leftarrow \max(\Delta, |v - V(s)|)$\;
    }
  }
  \tcp{Улучшение политики}
  $\textsc{политика-стабильна} \leftarrow \mathtt{true}$\;
  \For{каждого $s \in \cS$}{
    $a_{\text{старое}} \leftarrow \pi(s)$\;
    $\pi(s) \leftarrow \displaystyle\argmax_{a} \sum_{s',r} p(s',r \mid s,a)\bigl[r + \gamma\,V(s')\bigr]$\;
    \If{$a_{\text{старое}} \neq \pi(s)$}{$\textsc{политика-стабильна} \leftarrow \mathtt{false}$}
  }
}
\KwOut{$V \approx v^*$, $\pi \approx \pi^*$}
\end{algorithm}

\subsection*{Конечная сходимость}

\begin{theorem}[Конечная сходимость итерации политики]
\label{thm:pi-convergence}
\index{policy iteration!convergence}
Итерация политики завершается за конечное число внешних итераций.
\end{theorem}

\begin{proof}
Каждая политика $\pi$ — детерминированное отображение из $\cS$ в $\cA$.  Таких отображений не более $|\cA|^{|\cS|}$\,---\,конечное число.  Теорема об улучшении политики гарантирует, что всякий раз, когда порождается $\pi' \ne \pi$, $v^{\pi'} \ge v^\pi$ с хотя бы одним строгим неравенством.  Следовательно, последовательность политик, порождаемых итерацией политики, строго улучшается по ценности (пока не стабилизируется), что исключает двукратное посещение любой политики.  Поскольку политик конечное число, алгоритм должен завершиться.
\end{proof}

\begin{remark}
На практике итерация политики сходится за очень мало итераций\,---\,нередко менее десяти\,---\,даже для больших MDP.  Каждая итерация требует решения линейной системы (или итеративной оценки до сходимости), что доминирует в вычислениях.
\end{remark}

\begin{example}[Итерация политики на решётке]
\label{ex:pi-gridworld}
Возвращаясь к $4\times4$ решётке из примера~\ref{ex:gridworld-eval}, запускаем итерацию политики с равновероятной случайной политики.

\begin{itemize}[leftmargin=*]
\item \textbf{Итерация 0}: Оцениваем случайную политику (как вычислено выше).  Жадное улучшение немедленно даёт значительно лучшую политику: в каждом состоянии двигаться к ближайшему терминалу.
\item \textbf{Итерация 1}: Улучшенная политика уже выглядит почти оптимально.  Повторная оценка показывает, что она достигает $v^{\pi_1}(s) = -$ (манхэттенское расстояние до ближайшего терминала) для всех $s$.
\item \textbf{Итерация 2}: Жадная политика относительно $v^{\pi_1}$ совпадает с $\pi_1$.  Итерация политики останавливается.
\end{itemize}
Таким образом, итерация политики сходится всего за 2 внешние итерации на этой небольшой задаче, несмотря на то что итеративная оценка на каждом шаге требует многих внутренних проходов.
\end{example}

% ===========================================================
\section{Итерация ценности}
\label{sec:value-iteration}
\index{value iteration}

Итерация политики может быть дорогостоящей, если оценка политики требует многих проходов до сходимости перед каждым шагом улучшения.  \emph{Итерация ценности} избегает этого, чередуя один шаг оценки с шагом улучшения, применяя по существу \emph{оператор оптимальности Беллмана}\index{Bellman optimality operator} напрямую.

\subsection*{Оператор оптимальности Беллмана}

Определим $T^*$ как
\begin{equation}
\label{eq:bellman-optimality-op}
(T^* v)(s) = \max_{a \in \cA} \sum_{s' \in \cS}\sum_r p(s',r \mid s,a)\bigl[r + \gamma\, v(s')\bigr].
\end{equation}
Оптимальная функция ценности $v^*$ является единственной неподвижной точкой $T^*$: $T^* v^* = v^*$.

\begin{theorem}[Сжатие $T^*$]
\label{thm:tstar-contraction}
$T^*$ является $\gamma$-сжатием на $(B(\cS), \|\cdot\|_\infty)$:
\[
\|T^* u - T^* v\|_\infty \le \gamma\,\|u - v\|_\infty.
\]
Следовательно, итераты $v_{k+1} = T^* v_k$ сходятся к $v^*$ из любой начальной точки $v_0$.
\end{theorem}

\begin{proof}
Для любого $s \in \cS$,
\begin{align*}
(T^* u)(s) - (T^* v)(s)
&= \max_a \sum_{s',r} p(s',r\mid s,a)\,\gamma\, u(s')
- \max_a \sum_{s',r} p(s',r\mid s,a)\,\gamma\, v(s') \\
&\le \max_a \sum_{s',r} p(s',r\mid s,a)\,\gamma\bigl[u(s') - v(s')\bigr] \\
&\le \gamma\,\|u - v\|_\infty,
\end{align*}
используя $\max f - \max g \le \max(f - g)$ и неотрицательность весов переходов.  Симметрично, $(T^* v)(s) - (T^* u)(s) \le \gamma\|u-v\|_\infty$.  Взяв супремум-норму, получаем результат.
\end{proof}

\subsection*{Алгоритм}

\begin{algorithm}[ht]
\caption{Итерация ценности}
\label{alg:value-iteration}
\KwIn{Динамика MDP $p$, дисконтирование $\gamma$, порог $\theta > 0$}
Инициализировать $V(s) \leftarrow 0$ для всех $s \in \cS$; $V(\text{терминал}) \leftarrow 0$\;
\Repeat{$\Delta < \theta$}{
  $\Delta \leftarrow 0$\;
  \For{каждого $s \in \cS$}{
    $v \leftarrow V(s)$\;
    $V(s) \leftarrow \displaystyle\max_a \sum_{s',r} p(s',r \mid s,a)\bigl[r + \gamma\, V(s')\bigr]$\;
    $\Delta \leftarrow \max(\Delta,\; |v - V(s)|)$\;
  }
}
Вывести детерминированную жадную политику $\pi$ относительно $V$\;
\KwOut{$V \approx v^*$, $\pi \approx \pi^*$}
\end{algorithm}

\paragraph{Скорость сходимости.}
Поскольку $T^*$ является $\gamma$-сжатием, ошибка удовлетворяет $\|v_k - v^*\|_\infty \le \gamma^k \|v_0 - v^*\|_\infty$.  Для желаемой точности $\varepsilon$ необходимое число итераций $k \ge \log(\varepsilon / \|v_0 - v^*\|_\infty) / \log \gamma$, которое растёт как $O\!\bigl((1-\gamma)^{-1}\log(1/\varepsilon)\bigr)$.  Меньшее $\gamma$ означает более быструю сходимость, но более сильное дисконтирование будущих вознаграждений.

\subsection*{Связь с итерацией политики}

Итерацию ценности можно рассматривать как \emph{итерацию политики с усечённой оценкой}: вместо выполнения оценки до сходимости выполняем ровно один проход $T^\pi$ перед улучшением.  Поскольку $T^{\pi_k} = T^*$, когда $\pi_k$ жадна относительно $v_k$ (по определению максимума в~\eqref{eq:bellman-optimality-op}), каждый проход итерации ценности одновременно является одним шагом и оценки, и улучшения.

Более общо, если оценка политики прерывается после $m$ проходов (а не $1$ или $\infty$), получаем семейство алгоритмов, интерполирующих между чистой итерацией ценности ($m=1$) и итерацией политики ($m \to \infty$).

\begin{example}[Итерация ценности на решётке]
\label{ex:vi-gridworld}
На той же $4\times4$ решётке применяем итерацию ценности, начиная с $V \equiv 0$.

После каждого прохода функция ценности быстро отражает структуру оптимального поведения.  Примерно после 10--15 проходов с порогом $\theta = 10^{-4}$ алгоритм сходится, и жадная политика, извлечённая из $V$, совпадает с оптимальной политикой, найденной итерацией политики.  Итерация ценности требует меньше суммарных проходов, чем итерация политики (которая нуждается во многих внутренних проходах на каждую внешнюю итерацию), когда число состояний мало, но преимущество меняется для больших MDP, поскольку итерация политики сходится за очень мало внешних итераций.
\end{example}

% ===========================================================
\section{Обобщённая итерация политики}
\label{sec:gpi}
\index{Generalized Policy Iteration (GPI)}

Оценка политики и улучшение политики — не жёстко последовательные шаги\,---\,их можно чередовать в любой пропорции.  Термин \emph{обобщённая итерация политики} (GPI) обозначает общую идею поддержания функции ценности $V$ и политики $\pi$ с многократным применением некоторой формы оценки (делающей $V$ более согласованным с $\pi$) и некоторой формы улучшения (делающей $\pi$ более жадной относительно $V$).

\paragraph{Два конкурирующих процесса.}\index{GPI!competing processes}
Два процесса тянут в противоположные стороны: оценка толкает $V$ к $v^\pi$ для \emph{текущей} $\pi$, тогда как улучшение изменяет $\pi$ на жадную относительно текущей $V$, что немедленно делает $V$ неточной оценкой для \emph{новой} $\pi$.  Несмотря на это противоречие, два процесса взаимодействуют в следующем смысле: каждый делает прогресс к одной и той же цели.  Взаимодействие можно визуализировать как две линии в пространстве значение-политика, притягиваемые друг к другу (см. рис.~\ref{fig:gpi}).

\begin{figure}[ht]
\centering
\begin{tikzpicture}[>=Stealth, thick,
  box/.style={draw, rounded corners=3pt, minimum width=2.8cm, minimum height=1cm,
               font=\small, align=center, fill=white},
  every node/.append style={font=\small}]

% Узлы
\node[box, fill=blue!10] (eval) at (0,0) {Оценка\\политики\\$V \to v^\pi$};
\node[box, fill=orange!15] (impr) at (5,0) {Улучшение\\политики\\$\pi \to \text{жадная}(V)$};

% Стрелки между ними
\draw[->, bend left=30] (eval) to node[above] {$\pi$ фиксировано, $V$ улучшается} (impr);
\draw[->, bend left=30] (impr) to node[below] {$V$ фиксировано, $\pi$ улучшается} (eval);

% Метка сходимости
\node[draw, dashed, rounded corners=3pt, fill=green!10, minimum width=2.2cm,
      minimum height=0.8cm, align=center] (opt) at (2.5, -2.2) {$v = v^*$\\$\pi = \pi^*$};
\draw[->, dashed, gray] (eval) -- (opt);
\draw[->, dashed, gray] (impr) -- (opt);

\node[above=0.7cm of eval, font=\normalsize\bfseries] {Обобщённая итерация политики};
\end{tikzpicture}
\caption{Рамка GPI.  Оценка политики и улучшение политики взаимодействуют многократно.  Каждый шаг движется к устойчивой неподвижной точке, где одновременно $V = v^*$ и $\pi = \pi^*$.  В итерации политики оценка выполняется до завершения перед каждым шагом улучшения.  В итерации ценности один проход выполняет обе функции.  Все модельно-свободные алгоритмы в последующих главах можно понять как приближённую GPI.}
\label{fig:gpi}
\end{figure}

\paragraph{GPI как универсальная рамка.}
Каждый крупный алгоритм RL вписывается в шаблон GPI.  В методе Монте-Карло (глава~\ref{ch:mc}) оценка выполняется усреднением выборочных доходов, а улучшение — $\varepsilon$-жадным шагом.  В временном разностном обучении (глава~\ref{ch:td}) оценка использует одношаговые оценки с бутстрэпингом.  В методах актор-критик (глава~\ref{ch:actor-critic}) актор выполняет улучшение, а критик — оценку.

\paragraph{GPI на собственных и чужих данных.}\index{on-policy}\index{off-policy}
Важное различие, ставшее критически важным в модельно-свободном RL, — используются ли данные оценки из той же самой оцениваемой политики (\emph{на собственных данных}) или из другой \emph{поведенческой политики} (\emph{на чужих данных}).  В чистом DP это различие несущественно, поскольку модель $p$ предоставляет точные математические ожидания; данные вообще не собираются.  Однако, как только мы переходим к методам на основе выборки в главах~\ref{ch:mc}--\ref{ch:td}, различие становится центральным.

% ===========================================================
\section{Асинхронное динамическое программирование}
\label{sec:async-dp}
\index{asynchronous DP}

Алгоритмы, описанные до сих пор, выполняют \emph{синхронные} проходы: каждое состояние обновляется ровно один раз за итерацию.  \emph{Асинхронное} DP ослабляет это требование, обновляя состояния в произвольном порядке, любое число раз, не дожидаясь полного прохода.

\paragraph{Обновления на месте.}\index{in-place DP}
В синхронной версии мы поддерживаем два массива $V$ (текущий) и $V'$ (следующий), записывая $V'(s) = T^\pi V(s)$, используя только значения из $V$.  В варианте \emph{на месте} мы поддерживаем один массив и немедленно перезаписываем $V(s)$ новым значением.  Поскольку обновление $V(s)$ может использовать уже обновлённые значения $V(s')$ при $s' \ne s$, версия на месте не является синхронной, однако на практике обычно сходится быстрее: обновление в состоянии $s'$ немедленно распространяется на любое состояние $s$, обновляемое позже в том же проходе.

\paragraph{Приоритетный обход.}\index{prioritized sweeping}
Не все состояния одинаково важно обновлять.  \emph{Приоритетный обход} поддерживает очередь с приоритетами и обновляет состояние $s$, у которого ошибка Беллмана $|v - T^\pi V(s)|$ наибольшая.  При обновлении состояния приоритеты всех его предшественников (состояний, из которых достижимо $s$) пересчитываются.  Это концентрирует вычисления на наиболее информативных обновлениях и может существенно сократить число обновлений, необходимых для сходимости.

\paragraph{DP в реальном времени.}\index{real-time DP}
В \emph{DP реального времени} (Barto, Bradtke и Singh, 1995) агент взаимодействует со средой в реальном времени и обновляет только фактически посещаемые состояния.  Это концептуальный мост между модельным DP и модельно-свободным обучением TD: если модель доступна, посещаемые состояния могут обновляться с помощью точных резервных вычислений Беллмана; если нет, резервные вычисления заменяются выборочными переходами.

\begin{remark}
Асинхронное DP сходится к $v^*$ (или $v^\pi$) при условии, что каждое состояние продолжает обновляться бесконечно часто\,---\,то есть ни одно состояние не игнорируется постоянно.  Это условие \emph{достаточных обновлений} легко выполнить, многократно посещая состояния в произвольном порядке.
\end{remark}

% ===========================================================
\section{Эффективность и ограничения динамического программирования}
\label{sec:dp-limitations}
\index{curse of dimensionality}

\paragraph{Полиномиальная сложность.}
Алгоритмы DP полиномиальны относительно $|\cS|$ и $|\cA|$.  Каждый проход оценки политики или итерации ценности требует $O(|\cS|^2 |\cA|)$ операций (для каждого из $|\cS|$ состояний обновление суммирует по $|\cA|$ действиям и $|\cS|$ следующим состояниям).  Итерация политики выполняет не более $|\cA|^{|\cS|}$ внешних итераций, но на практике сходится за очень мало (обычно $O(\log |\cS|)$).  Сравните это с полным перебором политик, требующим оценки всех $|\cA|^{|\cS|}$ политик.

\paragraph{Проклятие размерности.}\index{curse of dimensionality}
Несмотря на полиномиальную сложность в $|\cS|$, DP нередко неприменимо, поскольку само $|\cS|$ астрономически велико.  Робот с $n$ непрерывными углами суставов, дискретизированными по $d$ делений каждый, имеет $|\cS| = d^n$ состояний, что экспоненциально растёт с $n$.  Это \emph{проклятие размерности} Беллмана: пространство состояний взрывается с числом измерений.  У скромного робота с 10 суставами и 100 делениями на сустав $|\cS| = 10^{20}$\,---\,безнадёжно неприемлемо.

\paragraph{Сравнение с линейным программированием.}
Оптимальную функцию ценности $v^*$ можно также найти, решив линейную программу с $|\cS|$ переменными и $|\cS| \cdot |\cA|$ ограничениями.  Методы на основе ЛП могут в принципе использовать эффективные решатели ЛП, но они также плохо масштабируются с $|\cS|$ и в целом не конкурентоспособны с DP для больших задач.

\paragraph{Почему выборка необходима.}
Ограничения DP мотивируют модельно-свободные методы, изучаемые в остальной части книги.  В методах Монте-Карло (глава~\ref{ch:mc}) $v^\pi$ оценивается по выборочным траекториям без знания $p$.  В методах временных разностей (глава~\ref{ch:td}) бутстрэпинг позволяет делать инкрементные обновления из отдельных переходов.  Во всех случаях рамка GPI сохраняется, но точные резервные вычисления Беллмана заменяются приближёнными, выборочными.

Ключевое понимание таково: \emph{DP говорит нам, что вычислять; выборка говорит, как вычислять это, когда $p$ недоступно или пространство состояний слишком велико}.

% ===========================================================
\section{Числовой пример: аренда автомобилей Джека}
\label{sec:car-rental}
\index{Jack's Car Rental}

Чтобы проиллюстрировать итерацию политики на реалистичной задаче, опишем пример \emph{аренды автомобилей Джека}, классику из Саттона и Барто (2018).

\subsection*{Постановка задачи}

Джек управляет двумя пунктами аренды автомобилей, A и B.  Каждый день клиенты приезжают в каждый пункт, чтобы взять машину.  Аренда и возврат распределены примерно по закону Пуассона: в среднем 3 аренды и 3 возврата в день в пункте A и 4 аренды и 2 возврата в день в пункте B.  Джек получает \$10 за каждую арендованную машину (если она доступна) и может перемещать машины overnight между пунктами (до 5 машин в ночь по цене \$2 за машину).  Каждый пункт вмещает не более 20 машин.

MDP определяется следующим образом.  Состояние $(n_A, n_B)$ — число машин в каждом пункте в конце дня (до overnight-перемещений), от $0$ до $20$ на каждой площадке.  Действие $a \in \{-5,\ldots,5\}$ — чистое число машин, перемещаемых из A в B (отрицательное означает перемещение из B в A).  Вознаграждение — ожидаемый доход от аренды минус затраты на перемещение.  Пространство состояний имеет $21^2 = 441$ состояние, а пространство действий — $11$ действий, что делает задачу приемлемой для точного DP.

\subsection*{Эволюция политики}

На рисунке~\ref{fig:car-rental} схематично показано, как политика эволюционирует при итерации политики.

\begin{figure}[ht]
\centering
\begin{tikzpicture}[>=Stealth, scale=0.95,
  box/.style={draw, minimum width=3.5cm, minimum height=2.4cm, fill=blue!5, align=center, font=\small},
  arr/.style={->, thick}]

\node[box] (pi0) at (0,0) {$\pi_0$: Не перемещать\\машины нигде\\\emph{(ничего не делать)}};
\node[box] (pi1) at (4.5,0) {$\pi_1$: Перемещать\\из B в A\\\emph{(B имеет больше возвратов)}};
\node[box] (pi2) at (9,0) {$\pi_2$: Сбалансированная\\перевозка\\$\approx \pi^*$};

\draw[arr] (pi0) -- node[above, font=\scriptsize]{оценка + улучшение} (pi1);
\draw[arr] (pi1) -- node[above, font=\scriptsize]{оценка + улучшение} (pi2);

\node[below=0.3cm of pi0, font=\scriptsize, align=center] {Итерация 0\\\emph{(случайная инициализация)}};
\node[below=0.3cm of pi1, font=\scriptsize, align=center] {Итерация 1};
\node[below=0.3cm of pi2, font=\scriptsize, align=center] {Итерация 4--5\\\emph{(сходимость)}};
\end{tikzpicture}
\caption{Качественная эволюция политики при итерации политики для задачи аренды автомобилей.  Ранние политики грубы; после 4--5 итераций политика стабилизируется до оптимальной стратегии перевозки, перемещающей машины из B в A (поскольку B имеет более высокое ожидаемое число возвратов, хранение машин в A захватывает больше дохода от аренды).}
\label{fig:car-rental}
\end{figure}

\paragraph{Наблюдения.}
Оптимальная политика перемещает машины из пункта с более высоким ожидаемым числом возвратов (B) в пункт с более высоким ожидаемым числом аренды (A), тщательно балансируя затраты на перевозку \$2 за машину против ожидаемого прироста дохода от аренды.  Итерация политики сходится за 4--5 внешних итераций, каждая из которых требует решения линейной системы $441 \times 441$.  Это демонстрирует, как итерация политики эффективно обрабатывает структурированную динамику вознаграждений и переходов, достигая оптимума значительно быстрее, чем чистая итерация ценности или полный перебор.

% ===========================================================
\section{Резюме}
\label{sec:dp-summary}

\begin{itemize}[leftmargin=*]
\item Динамическое программирование предполагает известную модель $p(s',r \mid s,a)$ и точно решает задачу нахождения оптимальной функции ценности и политики.
\item \textbf{Итеративная оценка политики} многократно применяет оператор Беллмана $T^\pi$; сходимость гарантирована теоремой о сжимающем отображении со скоростью $\gamma$.
\item \textbf{Теорема об улучшении политики} гарантирует, что жадная политика относительно $v^\pi$ не хуже $\pi$; если она совпадает с $\pi$, то $\pi$ оптимальна.
\item \textbf{Итерация политики} чередует оценку и улучшение и завершается за конечное число итераций, поскольку последовательность политик строго улучшается и пространство политик конечно.
\item \textbf{Итерация ценности} применяет $T^*$ напрямую, объединяя оценку и улучшение в одном проходе, и геометрически сходится к $v^*$.
\item \textbf{Обобщённая итерация политики} — объединяющая рамка: все алгоритмы RL можно рассматривать как чередующиеся оценку и улучшение, отличающиеся лишь тем, насколько приближённым является каждый шаг.
\item \textbf{Асинхронное DP} обновляет состояния в гибком порядке; приоритетный обход концентрирует вычисления там, где ошибки Беллмана наибольшие.
\item DP полиномиально в $|\cS|$ и $|\cA|$, но непригодно для больших пространств состояний из-за проклятия размерности, мотивируя методы на основе выборки из глав~\ref{ch:mc} и~\ref{ch:td}.
\end{itemize}

% ===========================================================
\section*{Упражнения}
\addcontentsline{toc}{section}{Упражнения}

\begin{exercise}
\label{ex:dp-contraction-rate}
Пусть $v_0 \equiv 0$ и $\|v^*\|_\infty \le R_{\max}/(1-\gamma)$.
\begin{enumerate}[label=(\alph*)]
  \item Покажите, что $\|v_k - v^*\|_\infty \le \gamma^k R_{\max}/(1-\gamma)$ для итератов $v_{k+1} = T^* v_k$.
  \item При $\gamma = 0.99$ и $R_{\max} = 1$ сколько итераций итерации ценности необходимо для достижения $\|v_k - v^*\|_\infty \le 10^{-3}$?
\end{enumerate}
\end{exercise}

\begin{exercise}
\label{ex:dp-policy-improvement-stochastic}
Обобщите теорему об улучшении политики на \emph{стохастические} политики $\pi'$.  То есть предположим, что $\sum_a \pi'(a \mid s)\,q^\pi(s,a) \ge v^\pi(s)$ для всех $s$.  Докажите, что $v^{\pi'} \ge v^\pi$ поточечно.  \emph{Подсказка}: доказательство по существу то же, что и для детерминированных политик; оператор $T^{\pi'}$ по-прежнему является сжатием.
\end{exercise}

\begin{exercise}
\label{ex:dp-value-iter-policy}
После того как итерация ценности сошлась к $V \approx v^*$ с $\|V - v^*\|_\infty \le \varepsilon$, какую оценку можно дать субоптимальности жадной политики $\pi$, извлечённой из $V$?  То есть оцените $v^*(s) - v^\pi(s)$ через $\varepsilon$ и $\gamma$.  \emph{Подсказка}: дважды используйте аргумент сжатия.
\end{exercise}

\begin{exercise}
\label{ex:dp-gridworld-compute}
В $4\times4$ решётке с $\gamma=1$ и случайной политикой ($\pi(a \mid s) = 0.25$):
\begin{enumerate}[label=(\alph*)]
  \item Запишите уравнения Беллмана для четырёх угловых состояний и четырёх состояний в серединах рёбер.
  \item Проверьте численно (или алгебраически), что сошедшиеся значения в примере~\ref{ex:gridworld-eval} удовлетворяют этим уравнениям.
\end{enumerate}
\end{exercise}

\begin{exercise}
\label{ex:dp-number-of-policies}
Конечный MDP имеет $|\cS|=5$ состояний и $|\cA|=3$ действия.  Сколько детерминированных политик существует?  Итерация политики посещает каждую политику не более одного раза.  Предполагая, что каждая внешняя итерация требует $O(|\cS|^2|\cA|)$ работы, какова оценка суммарных вычислений в наихудшем случае?  Сравните с работой при полном переборе политик.
\end{exercise}

\begin{exercise}
\label{ex:dp-async}
Объясните, почему синхронный проход на месте (один массив, состояния обновляются слева направо в решётке) может сходиться быстрее, чем версия с двумя массивами.  Приведите интуитивный пример в $4\times4$ решётке, где обновление распространяется по нескольким состояниям за один проход.
\end{exercise}

\begin{exercise}[Усечённая оценка политики]
\label{ex:dp-truncated}
Определите \emph{$m$-шаговую итерацию политики} как итерацию политики, в которой фаза оценки выполняет ровно $m$ проходов $T^\pi$ перед каждым шагом улучшения.  Покажите, что:
\begin{enumerate}[label=(\alph*)]
  \item $m=\infty$ восстанавливает стандартную итерацию политики.
  \item $m=1$ восстанавливает итерацию ценности.
  \item Для всех конечных $m$ последовательность функций ценности не убывает и сходится к $v^*$.
\end{enumerate}
\end{exercise}

\begin{exercise}
\label{ex:dp-bellman-lp}
Покажите, что оптимальная функция ценности $v^*$ является \emph{наименьшим} супер-решением уравнения оптимальности Беллмана: то есть $v^*$ является покоординатным минимумом всех функций $v$, удовлетворяющих $v(s) \ge (T^* v)(s)$ для всех $s$.  Используйте это для формулировки нахождения $v^*$ как линейной программы.
\end{exercise}

\begin{exercise}[Модификация аренды автомобилей Джека]
\label{ex:dp-car-rental}
В задаче об аренде автомобилей из раздела~\ref{sec:car-rental} Джек нанимает второго работника.  Первый работник может перемещать одну машину в ночь бесплатно; любые дополнительные машины по-прежнему стоят \$2 каждая.  Кроме того, каждый пункт теперь вмещает до 25 машин, но взимается плата за парковку \$4 в ночь за каждый пункт, у которого более 10 машин.  Опишите, как бы вы изменили функцию вознаграждения и определение пространства состояний.  Какие изменения влияют на вероятности переходов?  Какие влияют только на вознаграждение?
\end{exercise}

\begin{exercise}
\label{ex:dp-on-off-policy}
В модельном DP мы утверждали, что различие на собственных/чужих данных несущественно.  Объясните точно, почему.  Затем объясните, что идёт не так, если мы пытаемся применить обновления \emph{на чужих данных} в модельно-свободной (основанной на выборке) версии оценки политики, и сформулируйте условие, необходимое для согласованности таких обновлений.
\end{exercise}
